\part*{I. PHƯƠNG PHÁP TỌA ĐỘ TRONG MẶT PHẲNG}
\addcontentsline{toc}{part}{I. PHƯƠNG PHÁP TỌA ĐỘ TRONG MẶT PHẲNG}
\section{Lý thuyết chung}
\subsection{Hệ tọa độ}
Trong mặt phẳng với hệ tọa độ $Oxy$ cho các điểm: $A\left(x_A;y_A\right), B\left(x_B;y_B\right), C\left(x_C;y_C\right)$.
\begin{itemize}
\item Tọa độ vectơ: $\overrightarrow{AB} = \left(x_B-x_A;y_B-y_A\right)$
\item Tọa độ trung điểm $J$ của đoạn thẳng $AB$, trọng tâm $G$ của tam giác $ABC$ lần lượt là: 
$$J\left( \dfrac{x_A+x_B}{2}; \dfrac{y_A+y_B}{2}\right); \quad G\left( \dfrac{x_A+x_B+x_C}{3}; \dfrac{y_A+y_B+y_C}{3}\right)$$
\end{itemize}
\subsection{Phương trình đường thẳng}
\subsubsection{ Vectơ chỉ phương và vectơ pháp tuyến của đường thẳng:}
\begin{itemize}
\item Vectơ $\overrightarrow u (\overrightarrow u \neq \overrightarrow 0) $ là \textbf{vectơ chỉ phương} của đường thẳng $d$ nếu nó có giá song song hoặc trùng với đường thẳng $d$.
\item Vectơ $\overrightarrow n (\overrightarrow n \neq \overrightarrow 0) $ là \textbf{vectơ pháp tuyến} của đường thẳng $d$ nếu nó có giá vuông góc với đường thẳng $d$.
\item Đường thẳng $ax + by + c = 0$ có một vectơ pháp tuyến là $\overrightarrow n = (a; b)$.
\item Hai đường thẳng song song có cùng vectơ chỉ phương (vectơ pháp tuyến).
\item Hai đường thẳng vuông góc có vectơ pháp tuyến của đường thẳng này là vectơ chỉ phương của đường thẳng kia.
\item Nếu $\overrightarrow u, \overrightarrow n$ lần lượt là vectơ chỉ phương, vectơ pháp tuyến của đường thẳng $d$ thì $\overrightarrow u. \overrightarrow n = 0$. Do đó, nếu $\overrightarrow u = (a;b)$ thì $\overrightarrow n = (b; -a)$.
\item Một đường thẳng có vô số vectơ pháp tuyến, vô số vectơ chỉ phương. Nếu $\overrightarrow n$ là một vectơ pháp tuyến (vectơ chỉ phương) của đường thẳng $d$ thì $k\overrightarrow n (k \neq 0)$ cũng là một vectơ pháp tuyến, vectơ chỉ phương của $d$.
\end{itemize} 
\subsubsection{Phương trình đường thẳng}
\begin{itemize}
\item \textbf{Phương trình tổng quát} của đường thẳng: 
\begin{equation}
ax + by + c = 0\quad (a^2+b^2>0)
\end{equation}
Đường thẳng đi qua điểm $M(x_0;y_0)$ và nhận $\overrightarrow n = (a; b)$ là vectơ pháp tuyến có phương trình dạng:
\begin{equation}
a(x-x_0) + b(y-y_0) = 0
\end{equation}
Đặc biệt: đường thẳng đi qua $(a; 0), (0; b)$ có \textbf{phương trình theo đoạn chắn}:
\begin{equation}
\dfrac{x}{a}+\dfrac{y}{b}=1
\end{equation}
* Đường thẳng đi qua $M(x_0; y_0)$ và nhận vectơ $\overrightarrow n=(p;q)$ làm vectơ chỉ phương, có \textbf{phương trình tham số} là:
\begin{equation}
\left\{\begin{array}{l} x=x_0+pt\\y=y_0+qt \end{array}\right.
\end{equation}
Có \textbf{phương trình chính tắc} là:
\begin{equation}
\dfrac{x-x_0}{p}=\dfrac{y-y_0}{q} \quad (p,q \neq 0)
\end{equation}
\textbf{Đặc biệt}: đường thẳng đi qua 2 điểm phân biệt $A\left(x_A;y_A\right), B\left(x_B;y_B\right)$ có phương trình dạng:
\begin{equation}
\dfrac{{x - x_A }}{{x_B - x_A }} = \dfrac{{y - y_A }}{{y_B - y_A }}
\end{equation}
\item Đường thẳng đi qua $M(x_0;y_0)$ và có hệ số góc $k$ thì có \textbf{phương trình đường thẳng với hệ số góc} dạng:
\begin{equation}
y=k(x-x_0)+y_0
\end{equation}
\textbf{Chú ý:}
\begin{itemize}
\item Không phải đường thẳng nào cũng có hệ số góc. Các đường thẳng dạng $x=a$ không có hệ số góc. Do vậy, khi giải các bài toán dùng hệ số góc, ta phải xét cả trường hợp đặc biệt này.
\item Nếu $\overrightarrow n = (a; b)$ là vectơ pháp tuyến của đường thẳng thì hệ số góc của nó là $k = -\dfrac{a}{b}, b\neq 0.$
\end{itemize}
\end{itemize}
\subsubsection{Vị trí tương đối của 2 điểm và 1 đường thẳng}
Cho $A\left(x_A;y_A\right),B\left(x_B;y_B\right)$ và đường thẳng $\Delta: ax + by+ c = 0$. Khi đó:
\begin{itemize}
\item Nếu $\left(ax_A + by_A+ c \right)\left(ax_B + by_B+ c \right) <0$ thì $A,B$ ở về hai phía khác nhau đối với $\Delta$.
\item Nếu $\left(ax_A + by_A+ c \right)\left(ax_B + by_B+ c \right) >0$ thì $A,B$ ở cùng một phía đối với $\Delta$
\end{itemize}
\subsection{Góc và khoảng cách}
\begin{itemize}
\item Góc giữa hai vectơ $\vec{v}, \vec{w}$ được tính dựa theo công thức:
\begin{equation} \label{goc2vt}
\cos (\vec{u}, \vec{w}) = \dfrac{\vec{u}.\vec{w}}{\left | \vec{v} \right|.\left | \vec{w} \right| }
\end{equation}
\item Giả sử $\overrightarrow n_1,\overrightarrow n_2$ lần lượt là vectơ pháp tuyến của các đường thẳng $d_1$ và $d_2$. Khi đó:
\begin{equation} \label{goc2dt}
cos(d_1, d_2) = \dfrac{{\left| {\overrightarrow n _1 .\overrightarrow n _2 } \right|}}{{\left| {\overrightarrow n _1 } \right|.\left| {\overrightarrow n _2 } \right|}}
\end{equation}
\item Độ dài vectơ $\vec{u}=(a;b)$ là: 
\begin{equation}
\left|\vec{u}\right| = \sqrt{a^2+b^2}
\end{equation}
\item Khoảng cách giữa hai điểm $A(x_A;y_A),B(x_B;y_B)$ là:
\begin{equation}\label{ab}
AB = \sqrt{\left(x_B-x_A\right)^2+\left(y_B-y_A\right)^2}
\end{equation}
\item Diện tích tam giác $ABC$ là:
\begin{equation} \label{dttg}
S=\dfrac{1}{2}\sqrt{\left ( AB.AC \right )^2 - \left (\overrightarrow{AB}.\overrightarrow{AC}  \right )^2}
\end{equation}
\item Khoảng cách từ điểm $M(x_0; y_0)$ đến đường thẳng $d:ax+by+c=0$ được tính bằng công thức:
\begin{equation}\label{kc}
d_{(M; d)} = \dfrac{{\left| {ax_0 + by_0 + c} \right|}}{{\sqrt {a^2 + b^2 } }}
\end{equation}
\end{itemize}
\subsection{Phương trình đường tròn}
\begin{itemize}
\item Đường tròn tâm $I(a; b)$, bán kính $R$ có dạng:
\begin{equation}
(x-a)^2 + (y-b)^2 = R^2
\end{equation}
\item Phương trình:
\begin{equation} \label{dtrtq}
x^2 + y^2 + 2ax + 2by + c = 0, \quad (a^2 + b^2-c > 0)
\end{equation}
cũng là phương trình đường tròn với tâm $I(- a ;-b)$ và bán kính $R=\sqrt{a^2+b^2-c}$.
\item Phương trình tiếp tuyến của đường tròn tại điểm $M(x_0 ; y_0)$
\begin{equation}
(x_0-a)(x-x_0) + (y_0-b)(y-y_0) = 0
\end{equation}
\item Vị trí tương đối của đường thẳng $\Delta$ và đường tròn $\left(C\right)$ tâm $I$, bán kính $R$.
\begin{itemize}
\item Nếu $d_{(I;\Delta)} > R$ thì $\Delta$ và $\left(C\right)$ không cắt nhau.
\item Nếu $d_{(I;\Delta)} = R$ thì $\Delta$ và $\left(C\right)$ tiếp xúc tại $I'$ là hình chiếu của $I$ lên $d$.
\item Nếu $d_{(I;\Delta)} < R$ thì $\Delta$ và $\left(C\right)$ cắt nhau tại hai điểm $M,N$. Khi đó trung điểm $H$ của $MN$ là hình chiếu của $I$ lên $MN$ và 
\begin{equation}
MN= 2\sqrt{R^2-d^2_{(I,\Delta)}}
\end{equation}
\end{itemize}
\end{itemize} 
\subsection{Phương trình Elip}
\begin{itemize}
\item Elip là tập hợp các điểm $M$ di động thỏa mãn $MF_1+MF_2=2a$ với $F_1, F_2$ cố định, $F_1F_2=2c$, $a>c>0$ là các số cho trước. 
\item $F_1(-c;0)$,$F_2(c;0)$ được gọi là \textbf{tiêu điểm}, $F_1F_2=2c$ được gọi là \textbf{tiêu cự}. $MF_1,MF_2$ là các \textbf{bán kính qua tiêu}.
\item Các điểm $A_1(-a;0)$, $A_2(a;0)$, $B_1(0;-b)$, $B_2(0;b)$ được gọi là các \textbf{đỉnh} của elip. Đoạn thẳng $A_1A_2=2a$ được gọi là \textbf{trục lớn}, $B_1B_2=2b$ được gọi là \textbf{trục nhỏ}.
\item \textbf{Phương trình chính tắc của Elip} có hai tiêu điểm $F_1(-c;0)$, $F_2(c;0)$ là:
\begin{equation} \label{pte}
\dfrac{x^2}{a^2}+\dfrac{y^2}{b^2}=1
\end{equation}
Trong đó $a>b>0, b^2=a^2-c^2$.
\item \textbf{Tâm sai} $e=\dfrac{c}{a}$.
\item Cho elip $(E)$ có phương trình chính tắc $\eqref{pte}$. Hình chữ nhật $PQRS$ với $P(-a;b)$, $Q(a;b)$, $R(a;-b)$, $S(-a;-b)$ được gọi là \textbf{hình chữ nhật cơ sở} của Elip.
\item Nếu $M\in (E)$ và $M,F_1,F_2$ không thẳng hàng thì đường thẳng phân giác ngoài của góc $\widehat{F_1MF_2}$ chính là tiếp tuyến của $(E)$ tại $M$.
\end{itemize}
\subsection{Tứ giác nội tiếp}
Một số lợi ích có thể sử dụng của tứ giác nội tiếp:
\begin{itemize}
\item Viết được phương trình đường tròn ngoại tiếp.
\item Sử dụng được tính chất: các góc nội tiếp chắn cùng 1 cung thì bằng nhau.
\item Chứng minh được 1 điểm cách đều các điểm khác.
\item ...
\end{itemize}
Các cách chứng minh tứ giác $ABCD$ nội tiếp:
\begin{enumerate}
\item Bốn đỉnh cùng cách đều 1 điểm.
\item Có hai góc đối diện bù nhau.
\item Hai đỉnh cùng nhìn đoạn thẳng (tạo bởi hai đỉnh còn lại) hai góc bằng nhau.
\item $MA.MB=MC.MD$, trong đó: $M=AB\cap CD$; hoặc $NA.ND=NC.NB$, với $N=AD \cap BC$.
\item $IA.IC=ID.IB$ với $I$ là giao điểm hai đường chéo.
\item Tứ giác đó là hình thang cân, hình chữ nhật, hình vuông, ...
\end{enumerate}
\section{Một số kĩ thuật cơ bản}
\subsection{Kĩ thuật xác định tọa độ điểm}
\subsubsection{Dựa vào hệ điểm} 
Xác định tọa độ điểm $M$ thỏa mãn điều kiện nào đó với hệ các điểm $A_1,A_2, ..., A_n$. Đối với bài toán này, ta đặt $M(x;y)$ và khai thác giả thiết.
\vidu{ Cho tam giác $ABC$ có trọng tâm $G(1;2)$, trực tâm $H(-1;3)$. Tìm tọa độ tâm đường tròn ngoại tiếp $I$ của tam giác.
}
\lgi
Giả sử $I(x;y)$. Ta có: $\overrightarrow{GH}=(-2;1); \overrightarrow{GI}=(x-1;y-2)$. Vì $\overrightarrow{GH}=-2\overrightarrow{GI}$ nên: \\
\cl{$\begin{cases}-2(x-1)=-2\\-2(y-2)=1\end{cases} \Leftrightarrow \begin{cases} x=2 \\y = \dfrac{3}{2} \end{cases}$}
Vậy $I \left(2; \dfrac{3}{2}\right).$
\subsubsection{Xác định tọa độ giao điểm của hai đường}
\subsubsection*{Giao của hai đường thẳng} 
Tọa độ giao điểm của hai đường thẳng $d_1:ax+by+c=0, d_2: mx+ny+p=0$ (nếu có) là nghiệm của hệ phương trình:
\begin{equation}
\begin{cases} ax+by+c=0 \\ mx+ny+p=0 \end{cases}
\end{equation}
\subsubsection*{Giao của đường thẳng và đường tròn} 
Cho đường thẳng $d: \begin{cases} x= x_0+mt \\y=y_0+nt\end{cases}$ và đường tròn $(C): (x-a)^2+(y-b)^2=R^2$. Tọa độ giao điểm (nếu có) của $d$ và $(C)$ là nghiệm của hệ phương trình:
\begin{equation}
\begin{cases} x= x_0+mt \\y=y_0+nt \\ (x-a)^2+(y-b)^2=R^2 \end{cases}
\end{equation}
\subsubsection*{Giao của đường thẳng và Elip} 
Cho đường thẳng $d: \begin{cases} x= x_0+mt \\y=y_0+nt\end{cases}$ và elip $\left(E\right): \dfrac{x^2}{a^2}+\dfrac{y^2}{b^2}=1$. Tọa độ giao điểm của $d$ và $\left(E\right)$ (nếu có) là nghiệm của hệ phương trình:
\begin{equation}
\begin{cases} x= x_0+mt \\y=y_0+nt \\ \dfrac{x^2}{a^2}+\dfrac{y^2}{b^2}=1 \end{cases}
\end{equation}
\subsubsection*{Giao của hai đường tròn} 
Tọa độ giao điểm của hai đường tròn:
$$\left(C_1\right):x^2+y^2+2a_1x+2b_1y+c_1=0; \quad \left(C_2\right):x^2+y^2+2a_2x+2b_2y+c_2=0 $$
(nếu có) là nghiệm của hệ phương trình:
\begin{equation}
\begin{cases}
x^2+y^2+2a_1x+2b_1y+c_1=0 \\ x^2+y^2+2a_2x+2b_2y+c_2=0
\end{cases}
\end{equation}
\vidu{ Cho hai đường tròn: $
\left(C_1\right): (x-1)^2+(y-2)^2=25;\quad \left(C_2\right):\left(x-\dfrac{7}{2}\right)^2+\left(y+\dfrac{1}{2} \right)^2=\dfrac{25}{2}.
$
Tìm tọa độ giao điểm (nếu có) của chúng.
}
\lgi
Tọa độ giao điểm (nếu có) của hai đường tròn là nghiệm của hệ phương trình:\\
\cl{$\begin{cases}x^2+y^2-2x-4y-20=0 \\ x^2+y^2-7x+y=0 \end{cases} \Leftrightarrow  \begin{cases}x-y=4 \\ x^2+y^2-7x+y=0 \end{cases}  \Leftrightarrow  \begin{cases}x-y=4 \\ \left[ \begin{array}{l}x=6 \\x=1 \end{array} \right. \end{cases}$}\\
Vậy hai đường tròn cắt nhau tại $A(6;2), B(1;-3).$
\subsubsection{Điểm thuộc đường} 
Để tìm tọa độ điểm $M$ thuộc đường thẳng $d: \begin{cases} x= x_0+mt \\y=y_0+nt\end{cases}$ thỏa mãn điều kiện nào đó. Ta lấy điểm $M(x_0+mt;y=y_0+nt)$ và áp dụng giả thiết, ta thu được phương trình ẩn $t$. Như thế, ta gọi là \textbf{tham số hóa} tọa độ điểm $M$.
\vidu{Cho điểm $A(2;-1)$. Tìm tọa độ điểm $M$ thuộc đường thẳng $d:2x-y-4=0$ sao cho $AM=\sqrt{2}$}
\lgi
Giả sử $M(m;2m-4)$. Ta có: $AM = \sqrt{(m-2)^2+(2m-3)^2}$. Khi đó:
$$AM=\sqrt{2}\Leftrightarrow 5m^2-16m+11 = 0 \Leftrightarrow \left[  \begin{array}{l}
m = 1 \\ m =\dfrac{11}{5}\end{array}\right.$$
Vậy các điểm cần tìm là $M_1(1;-2)$, $M_2\left(\dfrac{11}{5};\dfrac{2}{5}\right)$.
\subsection{Tìm tọa độ hình chiếu của một điểm lên một đường thẳng}
\begin{minipage}{0.42\textwidth}
\begin{center}
\definecolor{qqwuqq}{rgb}{0.12941176470588237,0.12941176470588237,0.12941176470588237}
\definecolor{xdxdff}{rgb}{0.6588235294117647,0.6588235294117647,0.6588235294117647}
\definecolor{uuuuuu}{rgb}{0.26666666666666666,0.26666666666666666,0.26666666666666666}
\definecolor{qqqqff}{rgb}{0.3333333333333333,0.3333333333333333,0.3333333333333333}
\begin{tikzpicture}[line cap=round,line join=round,>=triangle 45,x=1.0cm,y=1.0cm]
\clip(1.06,2.16) rectangle (7.1,5.4);
\draw[line width=0.4pt,color=qqwuqq,fill=qqwuqq,fill opacity=0.15] (3.3998994432424787,3.147674708967659) -- (3.3016283853904382,3.3356715152933014) -- (3.1136315790647964,3.2374004574412614) -- (3.2119026369168364,3.049403651115619) -- cycle; 
\draw (0.86,1.82)-- (7.9,5.5);
\draw [domain=1.06:7.1] plot(\x,{(-33.8336--7.04*\x)/-3.68});
\begin{normalsize}
\draw[color=black] (5.8,4.16) node {$d$};
\draw [fill=qqqqff] (2.6,4.22) circle (1.5pt);
\draw[color=qqqqff] (2.42,4.08) node {$M$};
\draw[color=black] (2.68,4.8) node {$\Delta$};
\draw [fill=uuuuuu] (3.2119026369168364,3.049403651115619) circle (1.5pt);
\draw[color=uuuuuu] (3.1,2.62) node {$H$};
\draw [fill=xdxdff] (7.661525354969574,5.37534279918864) circle (1.5pt);
\draw[color=xdxdff] (7.8,5.66) node {$C$};
\end{normalsize}
\end{tikzpicture}
\end{center}
\end{minipage}
\begin{minipage}{0.55 \textwidth}
Để tìm tọa độ hình chiếu $H$ của $M$ lên đường thẳng $d$ ta có 2 cách:
\begin{itemize}
\item Cách 1: Viết phương trình đường thẳng $\Delta$ đi qua $M$ và vuông góc với $d$. Điểm $H$ chính là giao điểm của $d$ và $\Delta$.
\item Cách 2: Tham số hóa tọa độ của $H \in d$ và dựa vào điều kiện $MH\perp d$.\\
\end{itemize}
\end{minipage}
\vidu{Tìm tọa độ hình chiếu $H$ của điểm $M(-1;-1)$ lên đường thẳng $d: x-y+2=0$.}
\lgi
\textbf{Cách 1}\\
 Đường thẳng $\Delta$ đi qua $M$ và vuông góc với đường thẳng $d$ có phương trình dạng:
$$1.(x+1)+1.(y+1) = 0 \Leftrightarrow x+y+2=0$$
Do $H = d\cap \Delta$ nên tọa độ của $H$ là nghiệm của hệ phương trình:
$$\begin{cases} x-y+2=0 \\x+y+2=0\end{cases}$$
Giải hệ ta được $H(-2;0)$.\\
\textbf{cách 2}\\
Đường thẳng $d$ có vectơ chỉ phương $\vec{u}=(1;1)$. Giả sử $H(h;h+2) \in d$. Ta có: $\overrightarrow {MH} =(h+1;h+3)$
$$\overrightarrow {MH}. \vec{u} =0 \Leftrightarrow 1.(h+1) + 1.(h+3) = 0 \Leftrightarrow h= -2$$
Vậy $H(-2;0)$.
\subsection{Tìm tọa độ điểm đối xứng của một điểm qua một đường thẳng}
\begin{minipage}{0.55 \textwidth}
Để tìm tọa độ điểm đối xứng $M'$ của $M$ qua đường thẳng $d$ ta có 2 cách:
\begin{itemize}
\item Cách 1: Tìm tọa độ điểm $H$ là hình chiếu của $M$ lên $d$. Do $H$ là trung điểm $MM'$ nên áp dụng công thức tìm tọa độ trung điểm, ta tìm được $M$ 
\item Cách 2: Giả sử $M'(x;y)$ và $H$ là trung điểm của $MM'$. Khi đó ta có:
$$\begin{cases}H \in d \\ \overrightarrow {MM'}. \vec{u} =0 \end{cases}  $$
\end{itemize}
\end{minipage}
\begin{minipage}{0.43 \textwidth}
\begin{center}
\definecolor{qqwuqq}{rgb}{0.12941176470588237,0.12941176470588237,0.12941176470588237}
\definecolor{uuuuuu}{rgb}{0.26666666666666666,0.26666666666666666,0.26666666666666666}
\definecolor{qqqqff}{rgb}{0.3333333333333333,0.3333333333333333,0.3333333333333333}
\begin{tikzpicture}[line cap=round,line join=round,>=triangle 45,x=1.0cm,y=1.0cm]
\clip(1.32,1.7) rectangle (7.26,5.26);
\draw[line width=0.4pt,color=qqwuqq,fill=qqwuqq,fill opacity=0.15] (3.3998994432424787,3.147674708967659) -- (3.3016283853904382,3.3356715152933014) -- (3.1136315790647964,3.2374004574412614) -- (3.2119026369168364,3.049403651115619) -- cycle; 
\draw (0.86,1.82)-- (7.9,5.5);
\draw [domain=1.32:7.26] plot(\x,{(-33.8336--7.04*\x)/-3.68});
\begin{normalsize}
\draw[color=black] (5.8,4.16) node {$d$};
\draw [fill=qqqqff] (2.6,4.22) circle (1.5pt);
\draw[color=qqqqff] (2.42,4.08) node {$M$};
\draw[color=black] (2.68,4.8) node {$\Delta$};
\draw [fill=uuuuuu] (3.2119026369168364,3.049403651115619) circle (1.5pt);
\draw[color=uuuuuu] (3.1,2.62) node {$H$};
\draw [fill=qqqqff] (3.8238052738336705,1.8788073022312384) circle (1.5pt);
\draw[color=qqqqff] (4.02,2.16) node {$M'$};
\end{normalsize}
\end{tikzpicture}
\end{center}
\end{minipage}
\vidu{Tìm tọa độ điểm $M'$ là đối xứng của điểm $M(1;1)$ qua đường thẳng $d: x+y+2=0$.}
\lgi
\textbf{Cách 1}\\
Đường thẳng $d$ có vectơ chỉ phương $\vec{u}=(1;-1)$. Hình chiếu của $M$ lên đường thẳng $d$ là $H(h;-h-2) \in d$. Ta có: $\overrightarrow {MH} =(h-1;-h-3)$. Do đó:
$$\overrightarrow {MH}. \vec{u} =0 \Leftrightarrow 1.(h-1) - 1.(-h-3) = 0 \Leftrightarrow h= -1$$
Vậy $H(-1;-1)$. \\
Do $H$ là trung điểm của $MM'$ nên: 
$$\begin{cases} x_{M'}=2x_H-x_M=-3 \\y_{M'}=2y_H-y_M=-3 \end{cases}$$
Vậy $M'(-3;-3)$.\\
\textbf{Cách 2}\\
Đường thẳng $d$ có vectơ chỉ phương $\vec{u}=(1;-1)$.\\
Giả sử $M'(x;y)$. Khi đó trung điểm $MM'$ là $H\left(\dfrac{x+1}{2};\dfrac{y+1}{2}\right) \in d$ và $\overrightarrow {MM'}. \vec{u} =0$. Ta có hệ:
$$\begin{cases} \dfrac{x+1}{2}+\dfrac{y+1}{2} + 2=0\\ 1.(x-1) - 1.(y-1) = 0 \end{cases} \Leftrightarrow \begin{cases} x=-3 \\y=-3 \end{cases}$$
Vậy $M'(-3;-3)$.
\subsection{Viết phương trình đường thẳng đi qua 1 điểm, cách 1 điểm cho trước một khoảng cho trước}
\begin{minipage}{0.43 \textwidth}
\begin{center}
\definecolor{qqwuqq}{rgb}{0.,0.39215686274509803,0.}
\definecolor{sqsqsq}{rgb}{0.12549019607843137,0.12549019607843137,0.12549019607843137}
\definecolor{uuuuuu}{rgb}{0.26666666666666666,0.26666666666666666,0.26666666666666666}
\begin{tikzpicture}[line cap=round,line join=round,>=triangle 45,x=1.0cm,y=1.0cm]
\clip(2.04,-2.22) rectangle (8.08,2.12);
\draw[color=qqwuqq,fill=qqwuqq,fill opacity=0.1] (6.010379530080103,1.5072739550020466) -- (6.12518078873751,1.3288904607805374) -- (6.303564282959019,1.4436917194379444) -- (6.188763024301612,1.6220752136594534) -- cycle; 
\draw[color=qqwuqq,fill=qqwuqq,fill opacity=0.1] (6.746728772336625,-1.580084064428584) -- (6.542904495970455,-1.5212991653148908) -- (6.484119596856762,-1.725123441681061) -- (6.687943873222932,-1.7839083407947542) -- cycle; 
\draw (1.12,-1.64)-- (7.18,2.26);
\draw (0.78,-0.08)-- (8.2,-2.22);
\draw (7.211796120715524,0.032439171539375145)-- (6.188763024301612,1.6220752136594534);
\draw (7.211796120715524,0.032439171539375145)-- (6.687943873222932,-1.7839083407947542);
\draw (6.8,1.2) node[anchor=north west] {$p$};
\begin{normalsize}
\draw[color=black] (4.74,1.02) node {$\Delta_1$};
\draw[color=black] (4.8,-1.62) node {$\Delta_2$};
\draw [fill=uuuuuu] (2.6886495618807627,-0.6304730542351527) circle (1.5pt);
\draw[color=uuuuuu] (2.68,-0.28) node {$M$};
\draw [fill=sqsqsq] (7.211796120715524,0.032439171539375145) circle (1.5pt);
\draw[color=sqsqsq] (7.36,0.32) node {$N$};
\draw [fill=uuuuuu] (6.188763024301612,1.6220752136594534) circle (1.5pt);
\draw [fill=uuuuuu] (6.687943873222932,-1.7839083407947542) circle (1.5pt);
\draw[color=black] (7.22,-0.84) node {$p$};
\end{normalsize}
\end{tikzpicture}
\end{center}
\end{minipage}
\begin{minipage}{0.55 \textwidth}
Để viết phương trình đường thẳng $\Delta$ đi qua điểm $M$ và cách điểm $N\left( x_N;y_N \right)$ một khoảng bằng $p$ ta thường giả sử vectơ pháp tuyến của đường thẳng là $\vec{n}=(a;b), \quad (a^2+b^2>0) $ và áp dụng công thức tính khoảng cách - công thức $\eqref{kc}$.
\end{minipage}
\vidu{ Viết phương trình đường thẳng $\Delta$ đi qua $A(1;3)$ và cách điểm $B(-2;1)$ một khoảng bằng $3$.}
\lgi
Giả sử $\vec{n}=(a;b), (a^2+b^2>0)$ là vectơ pháp tuyến của đường thẳng cần tìm. Phương trình đường thẳng có dạng:
$$a(x-1)+b(y-3)=0 \Leftrightarrow ax + by -a-3b=0$$
Khi đó:\\
\cl{$
d_{(B;\Delta)}=3 \Leftrightarrow  \dfrac{|-2a+b-a-3b|}{\sqrt{a^2+b^2}}=3  \Leftrightarrow  5a^2 -12ab=0 \Leftrightarrow  \left[ \begin{array}{l}b=0 \\ b=\dfrac{12}{5}a \end{array} \right.$}\\
\begin{itemize}
\item $b=0$, chọn $a=1$ ta có $\Delta_1: x-1=0$.
\item $b=\dfrac{12}{5}a$, chọn $a=5,b=12$ ta có $\Delta_2: 5x+12y-41=0$.
\end{itemize}
Vậy có 2 đường thẳng thỏa mãn yêu cầu bài toán là: $\Delta_1 : x-1=0; \Delta_2 : 5x+12y-41=0
$.
\subsection{Viết phương trình đường thẳng đi qua 1 điểm, tạo với 1 đường thẳng khác một góc cho trước}
\begin{minipage}{0.55 \textwidth}
Để viết phương trình đường thẳng $\Delta$ đi qua điểm $M$ và tạo với đường thẳng $d$ một góc bẳng $\alpha$ ta thường giả sử vectơ pháp tuyến của đường thẳng là $\vec{n}=(a;b), \quad (a^2+b^2>0) $ và áp dụng công thức tính góc - công thức $\eqref{goc2dt}$.\\
\end{minipage}
\begin{minipage}{0.43 \textwidth}
\begin{center}
\definecolor{qqwuqq}{rgb}{0.12941176470588237,0.12941176470588237,0.12941176470588237}
\definecolor{sqsqsq}{rgb}{0.12549019607843137,0.12549019607843137,0.12549019607843137}
\definecolor{uuuuuu}{rgb}{0.26666666666666666,0.26666666666666666,0.26666666666666666}
\begin{tikzpicture}[line cap=round,line join=round,>=triangle 45,x=1.0cm,y=1.0cm]
\clip(2.58,-2.38) rectangle (9.48,-0.54);
\draw [shift={(3.6523809523809523,-1.04)},color=qqwuqq,fill=qqwuqq,fill opacity=0.1] (0,0) -- (-19.463771796396152:0.6) arc (-19.463771796396152:0.:0.6) -- cycle;
\draw [shift={(8.400096865451745,-1.04)},color=qqwuqq,fill=qqwuqq,fill opacity=0.1] (0,0) -- (180.:0.6) arc (180.:199.46377179639614:0.6) -- cycle;
\draw (2.52,-1.04)-- (9.52,-1.04);
\draw [domain=2.58:9.48] plot(\x,{(--0.23644616525091477-0.3332107589771663*\x)/0.9428523691977767});
\draw (3.6523809523809523,-1.04)-- (6.026238908916349,-1.8789383505226385);
\draw (6.026238908916349,-1.8789383505226385)-- (8.400096865451745,-1.04);
\draw [domain=2.58:9.48] plot(\x,{(-9.515975683329286--0.8389383505226384*\x)/2.3738579565353968});
\draw [domain=2.58:9.48] plot(\x,{(-0.5953101768739666--0.8389383505226384*\x)/-2.3738579565353963});
\begin{normalsize}
\draw[color=black] (5.76,-0.78) node {$d$};
\draw [fill=uuuuuu] (3.6523809523809523,-1.04) circle (1.5pt);
\draw [fill=sqsqsq] (6.026238908916349,-1.8789383505226385) circle (1.5pt);
\draw[color=sqsqsq] (6.04,-2.12) node {$M$};
\draw [fill=uuuuuu] (8.400096865451745,-1.04) circle (1.5pt);
\draw[color=black] (4.8,-1.74) node {$\Delta_2$};
\draw[color=black] (7.8,-1.56) node {$\Delta_1$};
\end{normalsize}
\end{tikzpicture}
\end{center}
\end{minipage}
\vidu{
Viết phương trình đường thẳng $\Delta$ đi qua $M(2;1)$ và tạo với đường thẳng $d: 2x+3y+4=0$ một góc $45^o$.}
\lgi
Giả sử $\vec{n}=(a;b), (a^2+b^2>0)$ là vectơ pháp tuyến của đường thẳng cần tìm. Phương trình đường thẳng có dạng:
$$ax + by -2a-b=0$$
Khi đó:\\
\cl{$\cos (d;\Delta) = \dfrac{1}{\sqrt{2}}  \Leftrightarrow  \dfrac{|2a+3b|}{\sqrt{a^2+b^2}\sqrt{4+9}}=\dfrac{1}{\sqrt{2}}  \Leftrightarrow 5a^2 -24ab-5b^2=0  \Leftrightarrow  \left[ \begin{array}{l}a=5b \\ a=-\dfrac{1}{5}b \end{array} \right.
$}\\
\begin{itemize}
\item $a=5b$, chọn $b=1,a=5$ ta có $\Delta_1: 5x+y-11=0$.
\item $a=-\dfrac{1}{5}b$, chọn $b=5,a=-1$ ta có $\Delta_2: -x+5y-3=0$.
\end{itemize}
Vậy có 2 đường thẳng thỏa mãn yêu cầu bài toán là: $\Delta_1: 5x+y-11=0; \Delta_2: -x+5y-3=0.$
\subsection{Viết phương trình đường phân giác trong của một góc}
Để viết phương trình đường phân giác trong của góc $\widehat{BAC}$ ta có nhiều cách. Dưới đây là 3 cách thường sử dụng:\\
\begin{minipage}{0.55 \textwidth}
\textbf{Cách 1}:\\ Dựa vào tính chất đường phân giác là tập hợp các điểm cách đều hai đường thẳng $AB:ax+by+c=0$ và $AC: mx+ny+p=0$, ta có:
$$\dfrac{|ax+by+c|}{\sqrt{a^2+b^2}}=\dfrac{|mx+ny+p|}{\sqrt{m^2+n^2}}$$
Hai đường thu được là phân giác trong và phân giác ngoài của góc $\widehat{ABC}$.\\
\end{minipage}
\begin{minipage}{0.43 \textwidth}
\begin{center}
\definecolor{qqwuqq}{rgb}{0.12941176470588237,0.12941176470588237,0.12941176470588237}
\definecolor{xdxdff}{rgb}{0.6588235294117647,0.6588235294117647,0.6588235294117647}
\definecolor{qqqqff}{rgb}{0.3333333333333333,0.3333333333333333,0.3333333333333333}
\begin{tikzpicture}[line cap=round,line join=round,>=triangle 45,x=1.0cm,y=1.0cm]
\clip(2.16,-2.38) rectangle (10.02,1.26);
\draw [shift={(4.16,0.12)},color=qqwuqq,fill=qqwuqq,fill opacity=0.1] (0,0) -- (-74.27415519827318:0.3) arc (-74.27415519827318:-17.999484399432905:0.3) -- cycle;
\draw [shift={(4.16,0.12)},color=qqwuqq,fill=qqwuqq,fill opacity=0.1] (0,0) -- (-130.54882599711345:0.4) arc (-130.54882599711345:-74.27415519827319:0.4) -- cycle;
\draw[color=qqwuqq,fill=qqwuqq,fill opacity=0.1] (4.8484787231943915,-0.3267416164774759) -- (5.050228866676739,-0.3922922046125348) -- (5.115779454811798,-0.19054206113018698) -- (4.91402931132945,-0.12499147299512811) -- cycle; 
\draw[color=qqwuqq,fill=qqwuqq,fill opacity=0.1] (3.8057728788417435,-0.6203405775858981) -- (3.9436790257431666,-0.45915157471410456) -- (3.7824900228713734,-0.3212454278126812) -- (3.64458387596995,-0.48243443068447467) -- cycle; 
\draw[color=qqwuqq,fill=qqwuqq,fill opacity=0.1] (3.4578083928822605,-0.7007434368908649) -- (3.296619390010467,-0.5628372899894416) -- (3.1587132431090437,-0.7240262928612353) -- (3.3199022459808374,-0.8619324397626585) -- cycle; 
\draw[color=qqwuqq,fill=qqwuqq,fill opacity=0.1] (2.865426295631281,0.31757144662967207) -- (3.067176439113629,0.2520208584946132) -- (3.132727027248688,0.45377100197696096) -- (2.93097688376634,0.5193215901120198) -- cycle; 
\draw (4.16,0.12)-- (2.62,-1.68);
\draw (2.62,-1.68)-- (9.7,-1.68);
\draw (9.7,-1.68)-- (4.16,0.12);
\draw [domain=2.16:10.02] plot(\x,{(--4.0368136414120075-0.9625695868848455*\x)/0.27103466642541774});
\draw [domain=2.16:10.02] plot(\x,{(-1.0119958619035563--0.27103466642541774*\x)/0.9625695868848455});
\draw [domain=2.159999999999998:9.7] plot(\x,{(-8.1528--1.8*\x)/-5.54});
\draw (4.54703167089238,-1.2545286552292212)-- (4.91402931132945,-0.12499147299512811);
\draw (4.634119859118297,-0.6952363802812701) -- (4.80531053259237,-0.7508578987385508);
\draw (4.655750449629462,-0.6286622294857975) -- (4.826941123103535,-0.6842837479430782);
\draw (2.6644073645623716,-0.30112015232686024)-- (2.93097688376634,0.5193215901120198);
\draw (2.7120967874273187,0.13691147812122034) -- (2.883287460901392,0.0812899596639397);
\draw (4.54703167089238,-1.2545286552292212)-- (3.64458387596995,-0.48243443068447467);
\draw (4.063893979075829,-0.9596215985298246) -- (4.180911225045395,-0.8228481941498123);
\draw (4.010704321816936,-0.9141148917638836) -- (4.1277215677865025,-0.7773414873838712);
\draw (2.6644073645623716,-0.30112015232686024)-- (3.3199022459808374,-0.8619324397626585);
\draw (3.0506634282563887,-0.5131395938547539) -- (2.933646182286822,-0.6499129982347663);
\draw (4.84,-1.9) node[anchor=north west] {$d$};
\draw (5.12,0.92) node[anchor=north west] {$d'$};
\begin{normalsize}
\draw [fill=qqqqff] (4.16,0.12) circle (1.5pt);
\draw[color=qqqqff] (4.3,0.48) node {$A$};
\draw [fill=qqqqff] (2.62,-1.68) circle (1.5pt);
\draw[color=qqqqff] (2.4,-1.84) node {$B$};
\draw [fill=qqqqff] (9.7,-1.68) circle (1.5pt);
\draw[color=qqqqff] (9.84,-1.94) node {$C$};
\draw[color=black] (3.18,4.44) node {$d$};
\draw[color=black] (-3.14,-1.58) node {$e$};
\draw [fill=xdxdff] (2.6644073645623716,-0.30112015232686024) circle (1.5pt);
\draw [fill=xdxdff] (4.54703167089238,-1.2545286552292212) circle (1.5pt);
\end{normalsize}
\end{tikzpicture}
\end{center}
\end{minipage}
\\
Sau đó, ta cần dựa vào vị trí tương đối của hai điểm $B,C$ với hai đường vừa tìm được để phân biệt phân giác trong, phân giác ngoài. Cụ thể, nếu $B,C$ ở cùng một phía thì đó là phân giác ngoài, ở khác phía thì là phân giác trong.\\
\\
\begin{minipage}{0.43 \textwidth}
\begin{center}
\definecolor{qqwuqq}{rgb}{0.12941176470588237,0.12941176470588237,0.12941176470588237}
\definecolor{uuuuuu}{rgb}{0.26666666666666666,0.26666666666666666,0.26666666666666666}
\definecolor{sqsqsq}{rgb}{0.12549019607843137,0.12549019607843137,0.12549019607843137}
\definecolor{qqqqff}{rgb}{0.3333333333333333,0.3333333333333333,0.3333333333333333}
\begin{tikzpicture}[line cap=round,line join=round,>=triangle 45,x=1.0cm,y=1.0cm]
\clip(0.04,-1.66) rectangle (7.54,1.8);
\draw [shift={(1.58,1.32)},color=qqwuqq,fill=qqwuqq,fill opacity=0.1] (0,0) -- (-100.61965527615514:0.4) arc (-100.61965527615514:-61.612197311387376:0.4) -- cycle;
\draw [shift={(1.58,1.32)},color=qqwuqq,fill=qqwuqq,fill opacity=0.1] (0,0) -- (-61.612197311387376:0.3) arc (-61.612197311387376:-22.60473934661962:0.3) -- cycle;
\draw (1.58,1.32)-- (1.16,-0.92);
\draw (1.16,-0.92)-- (6.96,-0.92);
\draw (1.58,1.32)-- (6.96,-0.92);
\draw [domain=1.58:7.540000000000002] plot(\x,{(--3.9707792092468477-1.7314256335367655*\x)/0.9357020517111803});
\draw (1.3466240092911175,0.07532804955262866)-- (2.5157020517111803,-0.4114256335367654);
\draw (2.5157020517111803,-0.4114256335367654)-- (2.749078042420062,0.8332463169106058);
\draw (3.2,-1.24) node[anchor=north west] {$d$};
\begin{normalsize}
\draw [fill=qqqqff] (1.58,1.32) circle (1.5pt);
\draw[color=qqqqff] (1.72,1.6) node {$A$};
\draw [fill=qqqqff] (1.16,-0.92) circle (1.5pt);
\draw[color=qqqqff] (0.8,-0.9) node {$B$};
\draw [fill=qqqqff] (6.96,-0.92) circle (1.5pt);
\draw[color=qqqqff] (7.24,-0.88) node {$C$};
\draw [fill=sqsqsq] (2.5157020517111803,-0.4114256335367654) circle (1.5pt);
\draw[color=sqsqsq] (2.38,-0.6) node {$D$};
\draw [fill=uuuuuu] (1.3466240092911175,0.07532804955262866) circle (1.5pt);
\draw[color=uuuuuu] (1.14,0.3) node {$B'$};
\draw [fill=uuuuuu] (2.749078042420062,0.8332463169106058) circle (1.5pt);
\draw[color=uuuuuu] (2.94,1.12) node {$C'$};
\end{normalsize}
\end{tikzpicture}
\end{center}
\end{minipage}
\begin{minipage}{0.55 \textwidth}
\textbf{Cách 2}: \\
 Lấy $B',C'$ lần lượt thuộc $AB,AC$ sao cho:
$$\overrightarrow{AB'}=\dfrac{1}{AB}.\overrightarrow{AB}; \overrightarrow{AC'}=\dfrac{1}{AC}.\overrightarrow{AC}.$$ Giả sử $\overrightarrow{AD}=\overrightarrow{AB'}+\overrightarrow{AC'}$ Khi đó tứ giác $AB'DC'$ là hình thoi (Vì sao?). 
\end{minipage}
\\Do đó, $\overrightarrow{AD}$ là vectơ chỉ phương của đường phân giác cần tìm.\\
\\
\textbf{Cách 3}: \\
Giả sử $\vec{u}=(a;b)$ là vectơ chỉ phương của đường phân giác cần tìm. Ta có: \\
\cl{$\cos (\overrightarrow{AB}, \vec{u}) = \cos (\overrightarrow{AC}, \vec{u}) \Leftrightarrow  \dfrac{\overrightarrow{AB}.\vec{u}}{\left|\overrightarrow{AB}\right|} = \dfrac{\overrightarrow{AC}.\vec{u}}{\left|\overrightarrow{AC}\right| }$}
\vidu{
Viết phương trình đường phân giác trong góc $A$ của tam giác $ABC$ biết $A(1;1)$, $B(4;5)$, $C(-4;-11)$.}
\lgi
\textbf{Cách 1.}
Ta có phương trình các cạnh:
$$AB: 4x-3y-1=0; AC: 12x-5y-7=0$$
Phương trình hai đường phân giác góc $A$ là:
$$\left[ \begin{array}{l}
\dfrac{4x-3y-1}{5}=\dfrac{12x-5y-7}{13}\\ \dfrac{4x-3y-1}{5}=-\dfrac{12x-5y-7}{13}
\end{array} \right. \Leftrightarrow \left[ \begin{array}{lc}
4x+7y-11=0 &  (d_1) \\56x-32y-24=0 & (d_2)
\end{array} \right.$$
Ta có:
$$\left(4x_C + 7y_C -11 \right)\left(4x_B + 7y_B-11 \right) <0$$
Do đó $B,C$ khác phía so với $(d_1)$ hay $(d_1)$ là đường phân giác cần tìm.\\
\textbf{Cách 2.}
Ta có $$\overrightarrow{AB} = (3;4); AB=5; \overrightarrow{AB'}=\dfrac{1}{5}\overrightarrow{AB}=\left(\dfrac{3}{5};\dfrac{4}{5}\right)$$
$$\overrightarrow{AC} = (-5;-12); AC=13; \overrightarrow{AC'}=\dfrac{1}{13}\overrightarrow{AC}=\left(-\dfrac{5}{13};-\dfrac{12}{13}\right)$$
Ta có: $\overrightarrow{AB'}+\overrightarrow{AC'}=\left(\dfrac{14}{65};-\dfrac{8}{65}\right)$. Vậy vectơ chỉ phương của đường phân giác cần tìm là: $\vec{u}= (7;-4)$.
Do đó phương trình đường phân giác cần tìm là:
\begin{displaymath}
4(x-1)+7(y-1)=0 \Leftrightarrow 4x+7y-11=0 \end{displaymath}
\textbf{Cách 3.}
Giả sử $\vec{u}=(a;b)$ là vectơ chỉ phương của đường phân giác cần tìm. Ta có $$\dfrac{\overrightarrow{AB}.\vec{u}}{\left|\overrightarrow{AB}\right|} = \dfrac{\overrightarrow{AC}.\vec{u}}{\left|\overrightarrow{AC}\right| } \Leftrightarrow \dfrac{3a+4b}{5}=\dfrac{-5a-12b}{13} \Leftrightarrow a=-\dfrac{7}{4}b$$
Vậy vectơ chỉ phương của đường phân giác cần tìm là: $\vec{u}= (7;-4)$.
Do đó phương trình đường phân giác cần tìm là:
\begin{displaymath}
4(x-1)+7(y-1)=0 \Leftrightarrow 4x+7y-11=0 \end{displaymath}
\subsection{Viết phương trình đường tròn đi qua ba điểm}
Để viết phương trình đường tròn đi qua ba điểm, ta sử dụng phương trình dạng $\eqref{dtrtq}$ và thay tọa độ ba điểm đó vào, thu được 1 hệ phương trình.\\
\vidu{Viết phương trình đường tròn ngoại tiếp tam giác $ABC$ biết:
$A(1;3)$,$B(-1;-1)$,$C(2;0)$.}
\lgi
Giả sử phương trình đường tròn $\left(C \right)$ cần tìm có dạng 
\begin{equation*}
x^2 + y^2 + 2ax + 2by + c = 0, \ \ \ (a^2 + b^2-c > 0)
\end{equation*}
Do $A,B,C \in \left(C \right)$ nên:
$$\begin{cases}1+9+2a+6b+c=0 \\ 1+ 1-2a-2b+c=0 \\4+2a+c=0\end{cases}\Leftrightarrow \begin{cases}a=0 \\ b=-1 \\c=-4\end{cases} \text{(Thỏa mãn)}$$
Vậy $\left(C\right):x^2+y^2-2y-4=0.$
\subsection{Viết phương trình đường thẳng đi qua hai tiếp điểm của đường tròn}
Cho điểm $A\left(x_A; y_A \right)$ nằm ngoài đường tròn $(C)$ tâm $I$ bán kính $R$. Từ $A$, kẻ hai tiếp tuyến $AT_1$, $AT_2$ tới $(C)$. Hãy viết phương trình đường thẳng $T_1,T_2$.\\
\\
Giả sử $T(x;y)$, $I(a;b)$ là tiếp điểm ($T$ là $T_1$ hoặc $T_2$). Khi đó, ta có:
\begin{equation} \label{pttd}
\begin{cases} T \in (C) \\ \overrightarrow{AT}.\overrightarrow{IT}=0
\end{cases} \Leftrightarrow \begin{cases} (x-a)^2+(y-b)^2=R^2\\ \left(x -x_A \right)(x-a)+\left(y -y_A \right)(y-b)=0 \end{cases}
\end{equation}
Trừ từng vế 2 phương trình của $\eqref{pttd}$ ta thu được 1 phương trình đường thẳng. Đó là phương trình cần tìm.\\
\vidu{ Cho đường tròn $(C)$ có phương trình $(x-4) ^{2}+y^{2}=4$ và điểm $M=(1; -2)$. Tìm tọa độ điểm $N$ thuộc $Oy$ sao cho từ $N$ kẻ được 2 tiếp tuyến $NA, NB$ đến $(C)$ ($A,B$ là tiếp điểm) đồng thời đường thẳng $AB$ đi qua $M$.}
\lgi
Gọi $I$ và $T$ lần lượt là tâm và tiếp điểm của đường tròn $(C)$ ($T$ là $A$ hoặc $B$). Ta có:
$$N\left( {0;n} \right),I\left( {4;0} \right),T\left( {x_0 ;y_0 } \right),\overrightarrow {NT} = \left( {x_0 ;y_0 - n} \right),\overrightarrow {IT} = \left( {x_0 - 4;y_0 } \right) $$
Khi đó
$$\begin{cases} T \in (C) \\ \overrightarrow {NT} .\overrightarrow {IT} = 0 \end{cases} \Leftrightarrow \begin{cases} x_0^2 + y_0^2 - 8x_0 + 12 = 0 \\ x_0^2 - 4x_0 + y_0^2 - ny_0 = 0 \end{cases}$$
Trừ từng vế hai phương trình của hệ, ta có: $4x_0 - ny_0 - 12 = 0 $.\\
Vậy $AB$ có phương trình là: $4x - ny - 12 = 0$. \\
Vì $AB$ đi qua $M(1;-2)$ nên: 
$$4 + 2n -12 = 0 \Rightarrow n = 4$$
 Vậy $N(0;4)$.
\section{Phương pháp giải toán}
Phương pháp chung để giải quyết bài toán hình học giải tích phẳng gồm các bước sau:
\begin{itemize}
\item Vẽ hình, xác định các yếu tố đã biết lên hình
\item Khám phá các tính chất khác của hình (nếu cần). Chú ý tìm các đường vuông góc, song song, đồng quy; các đoạn bằng nhau, góc bằng nhau; các góc đặc biệt; quan hệ thuộc giữa điểm và đường thẳng, đường tròn, ...
\item Xác định các điểm, đường thẳng (theo các kĩ thuật đã học) để thực hiện yêu cầu bài toán.
\end{itemize}
\vidu{Cho hình bình hành $ABCD$ có $A(-2;-1)$. Điểm $C$ thuộc đường thẳng $d: x-y-3=0$. Gọi $H,K,E$ lần lượt là hình chiếu của $A$ lên $CD,BD,BC$. Đường tròn ngoại tiếp tam giác $HEK$ là $(C):x^2+y^2+x+4y+3=0$. Tìm tọa độ các đỉnh $B$,$C$,$D$ }
\begin{minipage}{0.57 \textwidth}
\begin{center}
\definecolor{qqwuqq}{rgb}{0.,0.39215686274509803,0.}
\definecolor{uuuuuu}{rgb}{0.26666666666666666,0.26666666666666666,0.26666666666666666}
\definecolor{qqqqff}{rgb}{0.,0.,1.}
\begin{tikzpicture}[line cap=round,line join=round,>=triangle 45,x=1.0cm,y=1.0cm]
\clip(-3.88,-0.22) rectangle (5.46,3.76);
\draw [shift={(1.7222641509433962,0.1520754716981132)},color=qqwuqq,fill=qqwuqq,fill opacity=0.1] (0,0) -- (132.51044707800085:0.3) arc (132.51044707800085:161.32124082097388:0.3) -- cycle;
\draw [shift={(1.7222641509433962,0.1520754716981132)},color=qqwuqq,fill=qqwuqq,fill opacity=0.1] (0,0) -- (113.83168789897476:0.6) arc (113.83168789897476:132.51044707800085:0.6) -- cycle;
\draw [shift={(2.8,1.14)},color=qqwuqq,fill=qqwuqq,fill opacity=0.1] (0,0) -- (151.18920625702697:0.6) arc (151.18920625702697:180.:0.6) -- cycle;
\draw [shift={(-1.2,3.34)},color=qqwuqq,fill=qqwuqq,fill opacity=0.1] (0,0) -- (-28.810793742973072:0.3) arc (-28.810793742973072:0.:0.3) -- cycle;
\draw [shift={(-1.2,3.34)},color=qqwuqq,fill=qqwuqq,fill opacity=0.1] (0,0) -- (-47.48955292199915:0.6) arc (-47.48955292199915:-28.810793742973072:0.6) -- cycle;
\draw [shift={(-1.2,3.34)},color=qqwuqq,fill=qqwuqq,fill opacity=0.1] (0,0) -- (-137.48955292199915:0.6) arc (-137.48955292199915:-90.:0.6) -- cycle;
\draw [shift={(-0.823529411764706,1.8341176470588232)},color=qqwuqq,fill=qqwuqq,fill opacity=0.1] (0,0) -- (-165.96375653207355:0.3) arc (-165.96375653207355:-118.47420361007437:0.3) -- cycle;
\draw [shift={(-0.823529411764706,1.8341176470588232)},color=qqwuqq,fill=qqwuqq,fill opacity=0.1] (0,0) -- (-118.47420361007437:0.6) arc (-118.47420361007437:14.036243467926484:0.6) -- cycle;
\draw[color=qqwuqq,fill=qqwuqq,fill opacity=0.1] (-1.2,1.3521320343559642) -- (-1.412132034355964,1.3521320343559642) -- (-1.412132034355964,1.14) -- (-1.2,1.14) -- cycle; 
\draw[color=qqwuqq,fill=qqwuqq,fill opacity=0.1] (-0.8749789873074586,2.039915949229834) -- (-1.0807772894784693,1.9884663736870811) -- (-1.0293277139357166,1.7826680715160705) -- (-0.823529411764706,1.8341176470588232) -- cycle; 
\draw[color=qqwuqq,fill=qqwuqq,fill opacity=0.1] (1.8786381583381013,0.2954183118099261) -- (1.7352953182262882,0.4517923192046311) -- (1.5789213108315834,0.30844947909281817) -- (1.7222641509433962,0.1520754716981132) -- cycle; 
\draw (-1.2,3.34)-- (-3.6,1.14);
\draw (-3.6,1.14)-- (2.8,1.14);
\draw (-3.6,1.14)-- (5.2,3.34);
\draw (-1.2,3.34)-- (2.8,1.14);
\draw (-1.2,1.14)-- (1.7222641509433962,0.1520754716981132);
\draw (1.7222641509433962,0.1520754716981132)-- (0.8,2.24);
\draw (-0.823529411764706,1.8341176470588232)-- (-1.2,1.14);
\draw (-1.2,3.34)-- (5.2,3.34);
\draw [domain=-3.8800000000000003:5.2] plot(\x,{(--4.961569811320755-3.1879245283018873*\x)/-3.4777358490566037});
\draw (-1.2,3.34)-- (-1.2,1.14);
\draw (-1.2,3.34)-- (-0.823529411764706,1.8341176470588232);
\draw (-1.2,3.34)-- (1.7222641509433962,0.1520754716981132);
\begin{normalsize}
\draw [fill=qqqqff] (-1.2,3.34) circle (1.5pt);
\draw[color=qqqqff] (-1.34,3.48) node {$A$};
\draw [fill=qqqqff] (-3.6,1.14) circle (1.5pt);
\draw[color=qqqqff] (-3.72,1.46) node {$B$};
\draw [fill=qqqqff] (2.8,1.14) circle (1.5pt);
\draw[color=qqqqff] (3.14,1.12) node {$C$};
\draw [fill=qqqqff] (5.2,3.34) circle (1.5pt);
\draw[color=qqqqff] (5.14,3.6) node {$D$};
\draw [fill=uuuuuu] (1.7222641509433962,0.1520754716981132) circle (1.5pt);
\draw[color=uuuuuu] (1.9,0.) node {$H$};
\draw [fill=uuuuuu] (-1.2,1.14) circle (1.5pt);
\draw[color=uuuuuu] (-1.24,0.96) node {$E$};
\draw [fill=uuuuuu] (-0.823529411764706,1.8341176470588232) circle (1.5pt);
\draw[color=uuuuuu] (-0.68,2.12) node {$K$};
\draw [fill=uuuuuu] (0.8,2.24) circle (1.5pt);
\draw[color=uuuuuu] (0.94,2.52) node {$I$};
\draw[color=qqwuqq] (1.1,0.54) node {1};
\draw[color=qqwuqq] (1.22,0.9) node {2};
\draw[color=qqwuqq] (2.0,1.4) node {5};
\draw[color=qqwuqq] (-0.5,3.20) node {4};
\draw[color=qqwuqq] (-0.3,2.68) node {3};
\draw[color=qqwuqq] (-1.32,3.05) node {6};
\draw[color=qqwuqq] (-1.09,1.56) node {7};
\draw[color=qqwuqq] (-0.6,1.6) node {8};
\end{normalsize}
\end{tikzpicture}
\end{center}
\end{minipage}
\begin{minipage}{0.45 \textwidth}
\begin{center}
\textbf{Phân tích}
\end{center}
Bằng cách vẽ thử đường tròn ngoại tiếp tam giác $HKE$, ta dễ thấy nó đi qua tâm $I$ của hình bình hành. Ta cần chứng minh điều đó.\\
Tham số hóa điểm $C$ ta tìm được tọa độ của $I$ theo tham số. Vì $I \in (C)$ nên ta thay tọa độ của $I$ (có tham số $t$) vào phương trình $(C)$. Từ đó ta tìm được $t$. Dễ dàng suy ra tọa độ $A,I,C$.\\
Vì $\widehat{AHC}=\widehat{AEC}=90^o$ nên $H,E$ nằm trên đường tròn $(T)$ đường kính $AC$.
\end{minipage}
Do đó $H,E$ là giao điểm của $(C)$ và $(T)$. Ta tìm được $H,E$.\\
Viết phương trình $CH, CE$, tham số hóa rồi sử dụng tính chất trung điểm $I$, ta tìm được $B,D$.
\lgi
Gọi $I$ là tâm hình bình hành. \\
Ta có $\widehat{AHC}=\widehat{AEC}=90^o$ nên tứ giác $AEHC$ nội tiếp đường tròn $(T)$ đường kính $AC$. Suy ra:
\begin{equation} \label{gt1}
 \begin{cases} \widehat{H_1} = \widehat{C_5}= \widehat{A_4} \\ \widehat{H_2}=\widehat{A_3}  \end{cases} \Rightarrow \widehat{EHI} = \widehat{H_1}+\widehat{H_2} = \widehat{A_3}+\widehat{A_4} =\widehat{HAD}
\end{equation}
Mặt khác:
\begin{equation} \label{gt2}
\widehat{EAD} = \widehat{BAH} = 90^0  \Rightarrow \widehat{HAD} = \widehat{A_6}
\end{equation}
\begin{equation} \label{gt3}
\widehat{AEB} = \widehat{BKA} = 90^0  \Rightarrow AKEB \text{ nội tiếp } \Rightarrow \widehat{A_6}=\widehat{K_7} = 180^o-\widehat{K_8}
\end{equation}
Từ $\eqref{gt1}, \eqref{gt2}, \eqref{gt3}$ suy ra $\widehat{EHI}+\widehat{K_8}=180^o$. Do đó tứ giác $HEKI$ nội tiếp.\\
Giả sử $C(c;c-3) \in d$. Khi đó $I\left(\dfrac{c-2}{2};\dfrac{c-4}{2}\right)$. Vì $I \in (C)$ nên:
$$\left(\dfrac{c-2}{2}\right)^2+ \left(\dfrac{c-4}{2}\right)^2+\dfrac{c-2}{2}+2(c-4)+3=0 \Leftrightarrow \hay{c=2 \\c=-1}$$
Dễ thấy $\{ E, H \} = (C) \cap (T)$.
\begin{itemize}
\item \textbf{Trường hợp 1.} $C_1(2;-1), I_1(0;-1)$. Ta có phương trình đường tròn $(T):$
$$x^2+(y+1)^2 = 4$$
Tọa độ của $E,H$ là nghiệm của hệ phương trình:
$$\begin{cases}x^2+(y+1)^2 = 4 \\  x^2+y^2+x+4y+3=0 \end{cases} \Leftrightarrow \hay{(x;y) = \left(-\dfrac{8}{5};\dfrac{11}{5}\right) \\(x;y)=(0;-3)}$$
$C_1H, C_1E$ là hai đường thẳng có phương trình:
$$x-y-3=0; \quad  x-3y-5=0$$
$B_1,D_1$ thuộc hai đường thẳng trên và nhận $I_1$ làm trung điểm. Giả sử $(t;t-3),(3t'+5;t')$ là hai điểm đó. Ta có
$$\begin{cases}\dfrac{t+3t'+5}{2}= 0 \\ \dfrac{t+t'-3}{2} = -1\end{cases} \Leftrightarrow \begin{cases}t=4\\t'=-3\end{cases}$$
Vậy $B_1,D_1 \in \{ (4;1), (-4;-3) \}$
\item \textbf{Trường hợp 2.} $C_2(-1;-4)$. Bằng cách tương tự, ta có $B_2,D_2 \in \{(-1;0), (-2;-5) \}$
\end{itemize}
\vidu{Cho tam giác $ABC$ có điểm $A(2;3)$, tâm đường tròn ngoại tiếp $I(1;0)$, chân đường phân giác trong góc $A$ là $D(3;1)$. Tìm tọa độ các điểm $B$ và $C$.}
\lgi
\begin{minipage}{0.33 \textwidth}
\begin{center}
\definecolor{uuuuuu}{rgb}{0.26666666666666666,0.26666666666666666,0.26666666666666666}
\definecolor{xdxdff}{rgb}{0.49019607843137253,0.49019607843137253,1.}
\definecolor{qqqqff}{rgb}{0.,0.,1.}
\begin{tikzpicture}[line cap=round,line join=round,>=triangle 45,x=1.0cm,y=1.0cm]
\clip(-4.898756056057441,-2.7131091984349776) rectangle (-0.328360461492955,1.6211773986243208);
\draw(-2.706315364937946,-0.5543983641019417) circle (2.0018890225971977cm);
\draw (-4.38478643647689,-1.645404560602252)-- (-4.017413620360298,0.9584073152315373);
\draw (-4.017413620360298,0.9584073152315373)-- (-1.1464165124974368,-1.8090996149779999);
\draw [domain=-4.898756056057441:-0.328360461492955] plot(\x,{(-3.486113731516516-0.9454599622814062*\x)/0.32573833014068543});
\draw (-4.38478643647689,-1.645404560602252)-- (-1.1464165124974368,-1.8090996149779999);
\draw (-2.706315364937946,-0.5543983641019417)-- (-2.8073790106877983,-2.5537347035663425);
\begin{normalsize}
\draw [fill=qqqqff] (-2.706315364937946,-0.5543983641019417) circle (1.5pt);
\draw[color=qqqqff] (-2.571395937792131,-0.7567775048206639) node {$I$};
\draw [fill=xdxdff] (-4.017413620360298,0.9584073152315373) circle (1.5pt);
\draw[color=xdxdff] (-4.25788877711482,1.0646347616478349) node {$A$};
\draw [fill=xdxdff] (-4.38478643647689,-1.645404560602252) circle (1.5pt);
\draw[color=xdxdff] (-4.713241843731946,-1.6843485664481401) node {$C$};
\draw [fill=xdxdff] (-1.1464165124974368,-1.8090996149779999) circle (1.5pt);
\draw[color=xdxdff] (-0.9354978836491229,-1.7180784232345938) node {$B$};
\draw [fill=uuuuuu] (-2.8073790106877983,-2.5537347035663425) circle (1.5pt);
\draw[color=uuuuuu] (-2.589450436544719,-2.42521584539076) node {$D$};
\end{normalsize}
\end{tikzpicture}
\end{center}
\end{minipage}
\begin{minipage}{0.67 \textwidth}
Ta có:\\
Đường tròn tâm $I$ bán kính $IA$: $(x - 1)^2 + y^2 = 10$.\\
Đường thẳng $AD$: $2x + y - 7 = 0$.\\
\\
Gọi $E = AD \cap (I)$. Khi đó tọa độ điểm $E$ là nghiệm của hệ phương trình:
$$\va{ 2x + y - 7 = 0 \\  (x - 1)^2  + y^2  = 10 } \Leftrightarrow \va{ x = 4 \\  y =  - 1 }$$
Vậy $E(4;-1).$
\end{minipage}
$$\widehat{CAE} = \widehat{EAB} \Rightarrow EC = EB \Rightarrow IE \bot BC$$

Đường thẳng $BC$ đi qua điểm $D$ và có vectơ pháp tuyến $\vt{IB}$ nên có phương trình: $3x-y-8=0$.\\
Tọa độ của $B$ và $C$ là nghiệm của hệ:
$$
\va{
 3x - y - 8 = 0 \\ 
 (x - 1)^2  + y^2  = 10 } \Leftrightarrow \va{
 x = \frac{{5 - \sqrt 3 }}{2};y =  - \frac{{1 + 3\sqrt 3 }}{2} \\ 
 x = \frac{{5 + \sqrt 3 }}{2};y = \frac{{3\sqrt 3  - 1}}{2} }
$$
Vậy $B,C \in \left\{ {\left( {\frac{{5 - \sqrt 3 }}{2}; - \frac{{1 + 3\sqrt 3 }}{2}} \right),\left( {\frac{{5 + \sqrt 3 }}{2};\frac{{3\sqrt 3  - 1}}{2}} \right)} \right\}$